\documentclass[10pt,a4paper]{amsart}
\usepackage[margin=19mm]{geometry}
\usepackage{amsmath,amssymb,mathtools}
\usepackage[scr=rsfs,cal=euler]{mathalfa}
\usepackage{xfrac}
\usepackage[unicode,hidelinks,bookmarks=false]{hyperref}
\newcommand{\ie}{i.e.\ }
\newcommand{\cf}{cf.\ }
\newcommand{\resp}{respectively\ }
\newcommand{\Subsection}{\subsection}
\providecommand{\nobreakdash}{\nobreak\hbox{-}}
\setlength{\emergencystretch}{2em}
\multlinegap0pt
\usepackage{bm}
\usepackage{bbm}
\usepackage{dsfont}
\usepackage{xspace}
\usepackage{cancel}
\usepackage{mathtools}
\usepackage{amsthm}
\makeatletter
\@ifundefined{Theorem}{\newtheorem{Theorem}{Theorem}}{}
\@ifundefined{Proposition}{\newtheorem{Proposition}[Theorem]{Proposition}}{}
\makeatother
\newcommand\C{{\mathbb C}}
\renewcommand{\i}{\mathsf{i}}
\title[Geometry of the Space of Sections of Twistor Spaces with Cir]{Geometry of the Space of Sections of Twistor Spaces with Circle Action}
\author{Florian Beck, Indranil Biswas, Sebastian Heller, Markus R\"oser}
\date{Source: arXiv:2102.10853v3 (CC BY 4.0)}
\begin{document}
\maketitle

\noindent\emph{Source and licence.} Geometry of the Space of Sections of Twistor Spaces with Circle Action. Version arXiv:2102.10853v3 (CC BY 4.0), distributed under the Creative Commons Attribution 4.0 International licence, as recorded in the arXiv licence metadata for this exact version and shown on the abstract page https://arxiv.org/abs/2102.10853v3. Full rights notice: licenses/arxiv-2102.10853-CC-BY-4.0.txt.\par\smallskip

\noindent\textit{Editorial scope.} Counted unit: the Proposition whose source label is \texttt{Prop: energyC*-fixed\_sections} in the article (source0.tex lines 2655--2663, source label \texttt{Prop: energyC*-fixed\_sections}), delivered together with the article's own complete original proof of it (source0.tex lines 2665--2750, one \texttt{proof} environment of 86 lines). No mathematics is written, repaired or re-derived: statement and proof are the article's own text line for line. Required context retained uncounted: Eq: VHSblockmatrices (lines 2367--2382); eq:BB-slice2 (lines 2420--2424); eq:C*\_sectionLiftC (lines 2451--2453). Every definition, display and label of those lines is retained unchanged. Omitted from this package and not redistributed: 1--2366, 2383--2419, 2425--2450, 2454--2654, 2664--2664, 2751--3491. Nothing inside a retained excerpt is omitted or altered: every retained body is delivered exactly as the article's lines are written, so no omission receipt of any kind accompanies this package. Citations: the retained excerpts carry no citation command at all, so no bibliography entry of the article is transcribed and no reference is redistributed. Reference adaptation: none was needed and none was made; every reference inside the delivered unit resolves inside it, so no unresolved reference or citation is produced. Printed slips kept literal: no printed word, symbol, hypothesis or number of the counted statement or of its proof was altered. Layout and class: the article's own class is \texttt{sigma} with packages including ifpdf, tikz-cd, paralist, xstring; the wrapper is the lane's pinned \texttt{amsart} editorial wrapper at 10pt on A4, which restates the theorem-like environments the retained text needs and adds the article's own macro definitions that the retained text needs (\texttt{\textbackslash C}, \texttt{\textbackslash i}), with extra wrapper packages (bm, bbm, dsfont, xspace, cancel, mathtools). The delivered numbering is the unit's own. Assets: no retained span contains a figure environment, a graphics-inclusion command, a PDF-page inclusion, a table or a TikZ picture; no image file is copied into this package, the package declares no asset dependency and no image file is redistributed with it. Licence: version arXiv:2102.10853v3 of this article is distributed under the Creative Commons Attribution 4.0 International licence, as recorded in the arXiv licence metadata for this exact version and shown on the abstract page https://arxiv.org/abs/2102.10853v3. A full rights notice is delivered at licenses/arxiv-2102.10853-CC-BY-4.0.txt. Characters: every retained excerpt is pure ASCII, so the delivered engine, XeLaTeX with its default Latin Modern text font, reports no missing character.
\begin{equation}\label{Eq: VHSblockmatrices}
\overline{\partial} = \begin{pmatrix}
\overline{\partial}_{E_1} & 0 &\dots & \dots & 0\\
0 & \overline{\partial}_{E_2} & \ddots & & \vdots\\
\vdots & \ddots & \ddots & \ddots & \vdots\\
\vdots & & \ddots &\ddots & 0\\
0 & \dots & \dots & 0 & \overline{\partial}_{E_l}
\end{pmatrix},\qquad
\Phi = \begin{pmatrix}
0 & \dots & \dots & \dots & 0\\
\Phi^{(1)} & \ddots & & &\vdots\\
0 & \Phi^{(2)} & \ddots & &\vdots\\
\vdots & \ddots & \ddots & \ddots & \vdots\\
0 & \dots & 0 & \Phi^{(l-1)} & 0
\end{pmatrix},
\end{equation}

\begin{gather}
D''(\beta, \phi) + [\beta\wedge\phi] = {\overline{\partial}} \phi + [(\Phi+\phi)\wedge\beta]
 = 0,\nonumber \\
D'(\beta, \phi) = \partial^{h}\beta +[\Phi^{*_h}\wedge\phi] = 0.\label{eq:BB-slice2}
\end{gather}

\begin{equation}\label{eq:C*_sectionLiftC}
s_{(\beta,\phi)}(\lambda) = \left[\lambda, {\overline{\partial}} + \lambda(\Phi^{*_h}+\beta_1) + \sum_{j=2}^l\lambda^j\beta_j, \Phi + \lambda\partial^h + \sum_{j=0}^l\lambda^{j+1}\phi_j\right],
\end{equation}

\begin{Proposition}\label{Prop: energyC*-fixed_sections}
Let $\bigl({\overline{\partial}}, \Phi\bigr)$ be a stable $\C^*$-fixed $\operatorname{SL}(n,\C)$-Higgs bundle, and let $s_{(\beta,\phi)}$
be the $\C^*$-fixed section corresponding to \smash{$(\beta, \phi) \in \mathcal{B}^+_{({\overline{\partial}},\Phi)}$}.
Its energy is given by
\[
\mathcal{E}(s_{(\beta,\phi)}) = \mathcal{E}(s_0) = \sum_{k=2}^l(k-1)\deg(E_k),
\]
where $s_0$ is the twistor line through $\bigl({\overline{\partial}}, \Phi\bigr)$.
\end{Proposition}

\begin{proof}
Write $s_{(\beta,\phi)}$ in a form as in \eqref{eq:C*_sectionLiftC}. Then from the definition of
$\mathcal{E}$ it follows immediately that
\[
\mathcal{E}(s_{(\beta,\phi)}) = \mathcal{E}(s_0)+\frac{1}{2\pi \i} \int_\Sigma \operatorname{tr}(\Phi \wedge \beta_1).
\]
Next we will show that $\int_\Sigma \operatorname{tr}(\Phi \wedge \beta_1) = 0$. To this end, let us write
\[
\Phi = \sum_{k=1}^{l-1}\Phi^{(k)},\qquad \beta_1 = \sum_{k=1}^{l-1}\beta^{(k)},
\]
where $\Phi^{(k)} \in \Omega^{1,0}(\operatorname{Hom}(E_k, E_{k+1}))$, $ \beta^{(k)} \in
\Omega^{0,1}(\operatorname{Hom}(E_{k+1}, E_k))$; see the block form in \eqref{Eq: VHSblockmatrices}. It follows that
\[
\operatorname{tr}(\Phi\wedge\beta_1) = \sum_{k=1}^l\operatorname{tr}_{E_k}\bigl(\Phi^{(k-1)}\wedge\beta^{(k-1)}\bigr).
\]
Note that each of the above summands $\Phi^{(k-1)}\wedge\beta^{(k-1}$ belongs to
$\Omega^{1,1}(\mathrm{End}(E_k))$, and we have adopted the convention that $\Phi^{(k)}
= 0 = \beta^{(k)}$ if $k = 0, l$.

Now, equation \eqref{eq:BB-slice2} implies that
$
{\overline{\partial}} \phi_0 + [\Phi\wedge\beta_{1}] = 0
$
and we can write
\[
[\Phi\wedge\beta_1] = \sum_{k=1}^{l-1}\Phi^{(k-1)}\wedge\beta^{(k-1)} + \beta^{(k)}\wedge\Phi^{(k)}.
\]
Thus, for each $k = 1, \dots, l$,
\[
{\overline{\partial}} \phi_0^{(k)} + \Phi^{(k-1)}\wedge\beta^{(k-1)} + \beta^{(k)}\wedge\Phi^{(k)} = 0.
\]
Consider the case of $k = l$
\[
{\overline{\partial}} \phi_0^{(l)} + \Phi^{(l-1)}\wedge\beta^{(l-1)} = 0.
\]
Taking the trace of this equation and integrating over $\Sigma$, we find, using Stokes' theorem,
that
\[
\int_\Sigma\operatorname{tr}_{E_l}\bigl(\Phi^{(l-1)}\wedge\beta^{(l-1)}\bigr) = 0.
\]
Now assume that $\int_\Sigma\operatorname{tr}_{E_{k+1}}\bigl(\Phi^{(k)}\wedge\beta^{(k)}\bigr) = 0$
for all $k \geq k_0$. Then we have
\[
{\overline{\partial}} \phi_0^{(k_0)} + \Phi^{(k_0-1)}\wedge\beta^{(k_0-1)} + \beta^{(k_0)}\wedge\Phi^{(k_0)} = 0.
\]
Taking the trace and integrating yields
\begin{align*}
0 &= \int_\Sigma\operatorname{tr}_{E_{k_0}}\bigl(\Phi^{(k_0-1)}\wedge\beta^{(k_0-1)} + \beta^{(k_0)}\wedge\Phi^{(k_0)}\bigr) \\
&= \int_\Sigma\operatorname{tr}_{E_{k_0}}\bigl(\Phi^{(k_0-1)}\wedge\beta^{(k_0-1)}\bigr) - \int_\Sigma\operatorname{tr}_{E_{k_0+1}}\bigl(\Phi^{(k_0)}\wedge\beta^{(k_0)}\bigr) \\
&= \int_\Sigma\operatorname{tr}_{E_{k_0}}\bigl(\Phi^{(k_0-1)}\wedge\beta^{(k_0-1)}\bigr).
\end{align*}
It follows inductively that $\int_\Sigma\operatorname{tr}(\Phi\wedge\beta_1) = 0$.

It remains to compute the energy of the twistor line $s_0$. To this end, we observe that
\[
\mathcal{E}(s_0) = \frac{1}{2\pi \i}\int_\Sigma \operatorname{tr}(\Phi\wedge \Phi^{*_h}) =
 \frac{1}{2\pi \i}\int_\Sigma\sum_{k=2}^l \operatorname{tr}_{E_k}\bigl(\Phi^{(k-1)}\wedge \bigl(\Phi^{k-1}\bigr)^{*_h}\bigr)
 = \sum_{k=2}^l\mathcal E_k(s_0),
\]
where we put \smash{$\mathcal E_k(s_0) =
\frac{1}{2\pi \i}\int_\Sigma \operatorname{tr}_{E_k}\bigl(\Phi^{(k-1)}\wedge \bigl(\Phi^{k-1}\bigr)^{*_h}\bigr)$}
for $k \geq 2$.
The equation $F^{\nabla^h} + {[\Phi\wedge\Phi^{*_h}] = 0}$ is in the form of diagonal blocks, with respect
to the splitting $E = \bigoplus_{k=1}^lE_k$, with components
\[
F^{\nabla^h_{E_k}} + \Phi^{(k-1)}\wedge \bigl(\Phi^{(k-1)}\bigr)^{*_h} + \bigl(\Phi^{(k)}\bigr)^{*_h}\wedge \Phi^{(k)} = 0.
\]
This gives the following recursive relations:
\begin{align*}
\mathcal E_k(s_0) & = \frac{1}{2\pi \i}\int_\Sigma\operatorname{tr}_{E_k}\bigl(\Phi^{(k-1)}\wedge \bigl(\Phi^{(k-1)}\bigr)^{*_h}\bigr)\\
&= \frac{\i}{2\pi}\int_\Sigma\operatorname{tr}_{E_k}\bigl(F^{\nabla^h_{E_k}}\bigr) +\frac{1}{2\pi \i}\int_\Sigma\operatorname{tr}_{E_{k+1}}\bigl(\Phi^{(k)}\wedge\bigl(\Phi^{(k)}\bigr)^{*_h}\bigr)= \deg(E_k) + \mathcal E_{k+1}(s_0).
\end{align*}
Thus, if $k = l$, we find that
$
\mathcal E_l(s_0) = \deg(E_l) $,
and for general $k$ we get that
\[
\mathcal E_k(s_0) = \sum_{j=k}^{l-1}\deg(E_j) + \mathcal E_l(s_0) = \sum_{j=k}^l\deg(E_j).
\]
Therefore,
\[
\mathcal E(s_0) = \sum_{k=2}^l\mathcal E_k(s_0) = \sum_{k=2}^l\sum_{j=k}^l\deg(E_j)
 = \sum_{k=2}^l(k-1)\deg(E_k),
\]
and this completes the proof.
\end{proof}

\end{document}
